% !TeX root = ../imdt-topic-work-plan.tex
\begin{FormTable}
  \FormCell{53.3mm}{%
    论文题目：\field[117mm]{\ThesisTitle}\par
    或选题范围：\field[113mm]{\TopicScope}\par
    论文密级\textsuperscript{1}请在（\quad）中打“\tick”：公开\option{\PublicMark}\quad
    秘密\option{\SecretMark}\quad 机密\option{\ConfidentialMark}\par
    定密申请表编号：\field[43mm]{\ClassificationNumber}
    {\fontsize{9bp}{12bp}\selectfont （对于密级为秘密或机密的学位论文）}\par
    选题报告会日期：\ProposalMeetingDate\par
    论文选题是否涉及跨专业学位类别：是\textsuperscript{2}\option{\CrossYes}
    \field[31mm]{\CrossDisciplineName}\quad 否\option{\CrossNo}\par
    {\fontsize{9.5bp}{13bp}\selectfont
    若涉及，选题报告会是否已聘请所涉其他专业学位类别的专家参加：
    是\option{\ExpertYes}\quad 否\option{\ExpertNo}\par}
  }\\\hline
  \FormCell{53.7mm}{%
    \fontsize{10.5bp}{13.3bp}\selectfont
    选题报告要求：
    \begin{enumerate}
      \item 选题报告的字数一般不少于五千字；
      \item 选题报告的内容应包括：选题意义、该领域国内外研究动态（文献综述）、本课题研究的目的、预期成果、研究方案、研究方法及其论证、关键难点拟采取的解决措施、论文工作总体日程安排，预计答辩时间等；
      \item 要求查阅一定数量的中、外文献资料。文献综述部分不是将文献内容进行简单的堆砌，而应通过阅读、消化、提炼，对已有的研究成果和动态进行全面的综述；
      \item 选题报告的考核由两部分组成：书面报告和口头报告，分别评分后给出一个总成绩；
      \item 选题报告和研究工作应具有正确的政治立场，符合学术规范；
      \item 填好“研究生选题报告及论文工作计划”表格后，连同书面报告一起交院系研究生管理部门备案。
    \end{enumerate}
  }\\\hline
  \FormCell{82.8mm}{%
    阅读国内外文献情况：\par
    \hspace{1em}国内文献约\field[22mm]{\DomesticReferenceCount}篇，
    国外文献约\field[26mm]{\ForeignReferenceCount}篇。\par
    主要文献：\par
    \KeyReferences
  }\\\hline
\end{FormTable}
\vspace{2mm}
{\songti\fontsize{9bp}{12bp}\selectfont
\textsuperscript{1} 涉密学位论文的定密申请应当在论文开题之前进行。\par
\textsuperscript{2} 如选“是”，应按国家现行《研究生教育学科专业目录》规范填写所涉其他专业学位类别的名称。\par}
\clearpage
